Large language models can generate machine-readable recommendations for decomposing monolithic applications, but a valid output does not necessarily identify useful service boundaries. This paper presents TraceMicro, an evidence-linked workflow for assigning Java classes to candidate services, and examines its recorded behavior on a PetClinic development snapshot. The system combines source and static-dependency evidence, structured model responses, LangGraph orchestration, shared-class support, deterministic validation and bounded request accounting. We compare package grouping, a deterministic dependency baseline and three completed full-scope model executions over 25 eligible classes. The deterministic methods produce group sizes of 11/5/4/3/2 and 21/1/1/1/1, respectively. All three model executions satisfy coverage and citation-reference checks yet assign every class exclusively to one service. Together they consume 27 requests and 144,478 reported tokens. A separate three-class interactive pilot produces two overlapping groups, while three full-scope planning attempts stop before assignments begin. These observations demonstrate an inspectable execution pipeline and a repeatable single-group outcome. Compact v3 reaches its 1,024-token planning allowance. A 2,048-token variant completes planning and fails on the first assignment request. Against an AI-assisted author development reference prepared after model exposure, domain-pair F1 is 1.000 for package grouping and 0.703 for the dependency baseline and each original run. All five miss all three annotated shared classes. Independent architectural accuracy remains unmeasured.
